\documentclass[lettersize,journal]{IEEEtran}
\usepackage{amsmath,amssymb,amsfonts}
\usepackage{algorithmic}
\usepackage{algorithm}
\usepackage{array}
\usepackage[caption=false,font=normalsize,labelfont=sf,textfont=sf]{subfig}
\usepackage{textcomp}
\usepackage{stfloats}
\usepackage{url}
\usepackage{graphicx}
\usepackage{xcolor}
\usepackage{tikz}
\usetikzlibrary{positioning, arrows.meta, shapes.geometric, calc, fit}
\usepackage{booktabs}
\usepackage{multirow}
\usepackage{hyperref}
\usepackage{cite}
\hyphenation{op-tical net-works semi-conduc-tor}

\begin{document}

\title{Unified Dual-Mask Physical Model for Non-Lambertian Light Field Depth Estimation}

\author{Ruifeng Huang and Zhenglong Cui
\thanks{Manuscript received May 2026. This work was supported by Beihang University. (Corresponding author: Zhenglong Cui)}
\thanks{R. Huang and Z. Cui are with the School of Computer Science and Engineering, Beihang University, Beijing, China (e-mail: huangruifeng@buaa.edu.cn; czl@buaa.edu.cn).}}

\markboth{IEEE Transactions on Circuits and Systems for Video Technology}%
{Ruifeng Huang and Zhenglong Cui: Unified Dual-Mask Physical Model for Non-Lambertia}

\maketitle

\begin{abstract}

\end{abstract}

\begin{IEEEkeywords}

\end{IEEEkeywords}

\section{Introduction}

Light field imaging records the spatial and angular distribution of rays, providing dense multi-view sampling that supports robust depth estimation \cite{ng2005}. This high-dimensional representation underpins contemporary three-dimensional reconstruction, autonomous navigation, and virtual reality rendering \cite{ng2006}. While monocular depth estimation relies heavily on statistical priors \cite{eigen2014, godard2019}, the explicit angular gradients in light fields offer deterministic geometric constraints, yielding substantial reliability in structured environments \cite{honauer2016}.

The accuracy of conventional light field depth estimation relies on the photo-consistency assumption across angular views \cite{wanner2014}. Under ideal Lambertian conditions, surface points scatter incident light uniformly, generating linear traces in epipolar plane images (EPIs) that allow for precise disparity matching \cite{heber2018}. However, real-world environments are dominated by non-Lambertian surfaces, such as mirrors, glass, and polished metals, where the observed radiance is dictated by complex bidirectional reflectance distribution functions (BRDFs) \cite{blinn1977}. The observed pixel intensity $I(\mathbf{x}, \mathbf{u})$ becomes highly dependent on the angular coordinate $\mathbf{u} = (u, v)$, severely violating the photo-consistency criterion:
\begin{equation}
    I(\mathbf{x}, \mathbf{u}) \neq \mathcal{L}(\mathbf{x}), \quad \text{if} \quad \rho(\mathbf{x}) < 1,
\end{equation}
where $\mathbf{x}$ denotes the spatial coordinate and $\rho(\mathbf{x})$ represents the surface roughness relative to the incident wavelength. Consequently, the corresponding EPI traces are distorted by specular highlights and scattering artifacts, rendering standard EPI gradient operators mathematically unstable \cite{li2021}.

The performance degradation caused by non-Lambertian reflectance is severe and quantifiable. Existing EPI-based convolutional neural networks achieve a mean absolute error (MAE) of approximately 0.081 on standard Lambertian or uniformly textured mixed datasets \cite{lyu2022}. However, when evaluated on non-Lambertian regions, the estimation error escalates beyond 0.411. Visual inspection of these failure cases reveals that the network mechanically reproduces specular textures rather than inferring the underlying geometry. Current approaches typically mitigate this by treating specular pixels as statistical outliers \cite{zhang2023} or applying rigid data-driven priors \cite{peng2022}. These methods bypass the underlying physical light-matter interactions, lacking the interpretability required to decouple reflection components, and remain highly vulnerable to dynamic lighting variations.

Overcoming the non-Lambertian bottleneck requires transitioning from purely geometric correspondence matching to deterministic physical modeling. There is an urgent necessity for a framework that explicitly characterizes the dual nature of light-matter interactions—quantifying both the medium roughness and the angular wavevector deflection—to classify reflection types. By bridging this physical gap, it becomes possible to isolate true depth signals from specular interference, enabling accurate and robust depth reconstruction in complex, real-world scenarios.

Conventional light field depth estimation methods formulate disparity computation through explicit epipolar plane image (EPI) gradient extraction \cite{wanner2014, wang2015}. Methodologically, this framework relies on strict Lambertian photo-consistency, treating angular variations as linear geometric traces. At the technical level, hand-crafted gradient operators mathematically decompose when encountering specular highlights or scattering, as the EPI slopes are fractured by non-Lambertian radiance \cite{tao2013}. Furthermore, these approaches typically process isolated horizontal and vertical 2D EPI slices, discarding the vast majority of angular coordinates. This directional rigidity prevents the extraction of comprehensive angular frequency responses, leaving the underlying bidirectional reflectance distribution function (BRDF) completely unresolved.

Early deep learning approaches integrate convolutional neural networks with EPI analysis to learn implicit mappings from local angular patches to disparity \cite{heber2018, lyu2022}. Methodologically, these networks treat non-Lambertian disturbances as statistical outliers to be suppressed, rather than physical phenomena governed by light-matter interactions. Technically, restricting the input to 28 out of 81 available viewpoints limits the receptive field, preventing the network from capturing global contextual cues necessary to differentiate surface roughness. Without deterministic physical priors, data-driven EPI models struggle to distinguish genuine depth discontinuities from smooth artifacts, leading to persistent texture-copying issues in reflective regions.

Recent methods attempt to mitigate non-Lambertian interference by introducing robust loss functions, heuristic occlusion masks, or learned single-view priors \cite{li2021, zhang2023}. Methodologically, they lack an explicit mechanism to quantize surface roughness or model the Fresnel-specular versus cosine-diffuse reflectance components. At the technical level, because non-Lambertian pixels yield flat matching cost distributions, standard cost aggregation fails to identify a dominant disparity. Existing networks process these ambiguous regions blindly, lacking an uncertainty-aware module to propagate reliable depth from neighboring Lambertian areas. Without adaptive multi-modal energy optimization, the networks cannot dynamically attenuate specular energy or enhance diffuse signals, causing severe accuracy degradation on transparent and mirror surfaces \cite{wang2023}.

Therefore, a method is urgently required that integrates physical dual-mask modeling with full-viewpoint cost aggregation to resolve the inherent limitations of EPI-based frameworks in non-Lambertian scenes.

In this paper, we propose a unified dual-mask physical model and a multi-scale angular cost volume network (MACVO-Net) for non-Lambertian light field depth estimation. Departing from data-driven approaches that treat non-Lambertian effects as statistical outliers, our method explicitly formulates light-matter interactions as deterministic physical priors. By coupling an MRI-inspired angular frequency analysis with a full-viewpoint geometric cost volume, the proposed framework disentangles true depth discontinuities from specular and scattering artifacts, generating geometrically consistent dense disparity maps.

The main contributions of this work are summarized as follows:

\begin{enumerate}
    \item \textbf{Dual-Mask Physical Modeling and MRI-Inspired Angular Frequency Analysis.} We construct a dual-mask physical representation that combines a medium mask and an angular-direction mask to quantify pixel-level surface roughness and wavevector deflection. By introducing MRI k-space frequency analysis into light field imaging, we map the natural encoding of $9\times 9$ angular coordinates into distinct angular frequency responses. This establishes a deterministic frequency-domain mechanism that separates diffuse, Fresnel-specular, and scattering pixels via constrained least squares for BRDF parameter estimation. This formulation resolves the EPI slope fracture caused by violations of the photo-consistency assumption, providing deterministic physical interpretability for complex material identification. 
    
    \item \textbf{Adaptive Dual-Modal Energy Map Optimization with Geometric and Segmentation Priors.} We introduce a geometric consistency descriptor (GCD) to quantify surface roughness variations and design an adaptive dual-subtraction strategy for energy map refinement. By integrating scene ternary masks from the Segment Anything Model \cite{kirillov2023}, our method applies differentiated energy modulation—such as strongly suppressing smooth Lambertian regions by 85\% and amplifying smooth non-Lambertian regions by 50\%. This strategy isolates genuine depth signals from smooth artifacts, a task traditional EPI methods fail to accomplish, thereby yielding highly robust feature inputs by maximizing the non-Lambertian to Lambertian energy ratio.
    
    \item \textbf{Full-Viewpoint Cost Volume with Uncertainty-Aware Multi-Scale Network.} We develop MACVO-Net, which constructs an explicit $4D$ multi-disparity cost volume by warping and aggregating all 81 available viewpoints, discarding the directional rigidity of isolated 2D EPI slices. The aggregated cost volume is processed by a multi-scale U-Net \cite{ronneberger2015} architecture to capture comprehensive contextual information. We mathematically leverage the flat cost volume responses in non-Lambertian regions as explicit uncertainty signals, guiding the decoder to interpolate depth via spatial neighborhood propagation. This architecture resolves the texture-copying issue in reflective areas and surpasses the accuracy ceiling inherent to conventional EPI-based paradigms.
\end{enumerate}

Extensive experiments on standard light field benchmarks \cite{honauer2016} demonstrate the effectiveness of our approach. Quantitative evaluations indicate that the proposed framework achieves a Mean Absolute Error (MAE) of 0.14 and reduces the BadPix metric to 1.9\% on complex non-Lambertian scenes, outperforming existing state-of-the-art methods by a substantial margin.

\section{Related Work}

()

\section{Methodology}

In this section, we present the Multi-scale Angular Cost Volume Network (MACVO-Net), an end-to-end architecture designed for robust depth estimation in non-Lambertian scenarios. Different from existing frameworks that rely on implicit epipolar plane image (EPI) feature extraction \cite{wanner2014, heber2018}, MACVO-Net constructs an explicit 4D cost volume that aggregates photometric geometry from all available angular views. Coupled with a multi-scale U-Net structure \cite{ronneberger2015}, the model leverages the physical flatness of the cost volume in non-Lambertian regions as an explicit uncertainty signal, guiding the decoder to recover precise depth via spatial propagation. As illustrated in Fig.~\ref{fig:architecture}, the overall pipeline follows a two-branch design that converges into an encoder-decoder network for geometric regularization.

The input to MACVO-Net consists of a dense $9 \times 9$ angular array of light field images. We represent the raw light field data as a tensor $\mathbf{I} \in \mathbb{R}^{B \times 81 \times C_{in} \times H \times W}$, where $B$ denotes the batch size, $C_{in}$ represents the input color channels, and $H \times W$ corresponds to the spatial resolution of each sub-aperture image. Prior to feature extraction, the preprocessing module selects the center view $\mathbf{I}_{center} \in \mathbb{R}^{B \times C_{in} \times H \times W}$ from the 81 views to serve as the spatial alignment anchor. The remaining 80 non-center views are grouped for subsequent multi-view geometric computation, establishing the basis for cost aggregation.

The data flow bifurcates into two parallel feature extraction streams. First, the Feature Encoder processes $\mathbf{I}_{center}$ to capture spatial textures and structural information, projecting the input image into a 32-channel feature map $\mathbf{F}_{center} \in \mathbb{R}^{B \times 32 \times H \times W}$. Concurrently, the Cost Volume Builder evaluates the entire 81-view input across 32 predefined disparity hypotheses. For a given disparity $d$, the module applies spatial transformations to warp the 80 non-center views to the center perspective and computes the averaged photometric $L_1$ error:
\begin{equation}
Cost(b, d, h, w) = \frac{1}{80} \sum_{v=1}^{80} || \mathcal{W}(\mathbf{I}_v, d) - \mathbf{I}_{center} ||_1
\end{equation}
where $\mathcal{W}$ denotes the spatial warping operator. This yields a 4D cost volume of dimension $(B, 32, H, W)$. The Feature Fusion module then concatenates $\mathbf{F}_{center}$ and the cost volume along the channel dimension, producing a unified 64-channel representation that embeds both appearance cues and matching geometry. This fused tensor is subsequently fed into a U-Net Encoder comprising three sequential down-sampling stages. These stages progressively expand the channel dimension to 128 and 256, halving the spatial resolution at each step to aggregate high-level contextual information and expand the receptive field for handling occlusions.

The encoded multi-scale features are processed by the U-Net Decoder to regress the final depth map. Through three up-sampling stages, the decoder restores the spatial resolution while integrating high-frequency structural details from the encoder via skip connections. At the final decoder layer, a $1 \times 1$ convolution compresses the 32-channel output into a single-channel representation. A subsequent sigmoid activation function normalizes the values, producing the final dense disparity map $\mathbf{D} \in \mathbb{R}^{B \times 1 \times H \times W}$ with values strictly bounded between 0 and 1. This specific output format ensures direct compatibility with standard continuous depth estimation metrics and facilitates stable backpropagation during end-to-end training.

\subsection{}

The FeatureEncoder processes the center view $\mathbf{I}_{center}$ to extract high-level spatial textures and structural cues. Purely photometric 4D cost volumes are highly sensitive to occlusions and non-Lambertian reflections, where the constant intensity assumption fails \cite{li2021}. In these regions, cost volumes become ambiguous and lack the discriminative gradients required for direct depth regression. The FeatureEncoder mitigates this by decoupling semantic appearance from geometric matching. By processing the center view exclusively, this module establishes a robust spatial anchor that provides the subsequent U-Net regularization with surrounding structural boundaries, enabling depth propagation in regions where the cost volume signal is degraded.

Architecturally, the module maps the input center view tensor $\mathbf{I}_{center} \in \mathbb{R}^{B \times C_{in} \times H \times W}$ into a higher-dimensional feature space $\mathbf{F}_{center} \in \mathbb{R}^{B \times 32 \times H \times W}$. This transformation is executed without spatial downsampling to maintain strict pixel-wise alignment with the subsequent cost volume tensor. Let $C_{out} = 32$ denote the target channel dimension and $C_{hidden}$ denote the intermediate channel dimension. The feature encoding operator $\Phi(\cdot)$ is formulated as a composition of consecutive convolutional stages:

\begin{equation}
\mathbf{F}_{center} = \Phi(\mathbf{I}_{center}) = \sigma \left( \mathcal{F}_{3 \times 3}^{(C_{out}, C_{hidden})} \left( \sigma \left( \mathcal{F}_{3 \times 3}^{(C_{hidden}, C_{in})} \left( \mathbf{I}_{center} \right) \right) \right) \right)
\end{equation}

where $\mathcal{F}_{3 \times 3}^{(o, i)}(\cdot)$ denotes a parameterized convolutional block that maps $i$ input channels to $o$ output channels using a $3 \times 3$ kernel, a stride of one, and appropriate spatial padding to preserve the original spatial dimensions $H$ and $W$. The function $\sigma(\cdot)$ represents a Parametric Rectified Linear Unit (PReLU) activation function, which introduces non-linearity while preserving negative responses critical for texture representation.

The resulting feature map $\mathbf{F}_{center}$ serves as the structural prior for the network. Unlike existing light field architectures that rely on implicit epipolar plane image (EPI) feature extraction and concatenate spatial features early in the pipeline \cite{heber2018}, MACVO-Net isolates geometric matching from semantic appearance. The FeatureEncoder outputs $\mathbf{F}_{center}$ directly to the Feature Concatenation module, where it is fused with the explicit 4D cost volume along the channel dimension. This explicit fusion allows the high-frequency structural gradients within $\mathbf{F}_{center}$ to act as spatial anchors, effectively compensating for the flat photometric responses in non-Lambertian regions and guiding the downstream U-Net encoder.

\subsection{}

The CostVolumeBuilder translates the multi-view geometry of the light field into an explicit volumetric representation. Traditional approaches extract features from Epipolar Plane Images (EPIs) to implicitly infer disparity \cite{wanner2014}, \cite{heber2018}. These implicit formulations conflate geometric matching with local texture extraction, resulting in brittle performance when the Lambertian assumption is violated by specular or transparent materials \cite{li2021}, \cite{zhang2023}. In contrast, the CostVolumeBuilder computes a purely photometric 4D cost volume to isolate the geometric matching process. In non-Lambertian regions, photometric consistency fails, producing a flat cost profile across the disparity dimension. This flatness constitutes a critical uncertainty signal that guides the subsequent regularization network to prioritize semantic structural boundaries over unreliable photometric matching.

The module ingests the full light field tensor $\mathbf{I} \in \mathbb{R}^{B \times 81 \times C_{in} \times H \times W}$, comprising one center view $\mathbf{I}_c$ and 80 peripheral views. Let $\mathcal{D} = \{d_k\}_{k=1}^{32}$ denote the set of 32 predefined disparity hypotheses. For a given hypothesis $d_k$, each peripheral view $\mathbf{I}_v$ is warped to align with the center view coordinate system. Assuming the standard multi-view geometry of a plenoptic camera \cite{ng2005}, the disparity $d_k$ dictates a spatial shift proportional to the angular offset of view $v$. Let $\Delta \mathbf{p}_v = (u_v, v_v)$ represent the angular coordinate offset of view $v$ relative to the optical center. The warping operator $\mathcal{W}(\cdot, \cdot)$ synthesizes the warped image tensor $\hat{\mathbf{I}}_v(d_k)$ via differentiable bilinear sampling:

\begin{equation}
\hat{\mathbf{I}}_v(d_k) = \mathcal{W}(\mathbf{I}_v, \Delta \mathbf{p}_v \cdot d_k)
\end{equation}

where $\Delta \mathbf{p}_v \cdot d_k$ defines the 2D dense displacement field applied to the spatial grid of $\mathbf{I}_v$. This operation shifts the spatial coordinates of $\mathbf{I}_v$ such that points corresponding to the hypothesized depth $d_k$ align perfectly with their corresponding locations in $\mathbf{I}_c$.

The cost volume $\mathbf{C}$ is constructed by evaluating the photometric discrepancy between the warped peripheral views and the center view. To suppress sensor noise and mitigate the impact of moderate occlusions, the L1 photometric error is averaged across all 80 peripheral views. The construction of the cost volume element is formulated as:

\begin{equation}
\mathbf{C}(b, k, h, w) = \frac{1}{80} \sum_{v=1}^{80} \left\| \mathcal{W}(\mathbf{I}_v, \Delta \mathbf{p}_v \cdot d_k) - \mathbf{I}_c \right\|_1
\end{equation}

where $(b, k, h, w)$ index the batch, disparity hypothesis, height, and width dimensions, respectively. The L1 norm $\| \cdot \|_1$ calculates the absolute intensity difference across the $C_{in}$ channels, yielding a scalar penalty for each pixel. The output of this module is the 4D cost volume tensor $\mathbf{C} \in \mathbb{R}^{B \times 32 \times H \times W}$. The spatial resolution $(H, W)$ is strictly preserved to ensure pixel-wise alignment with the feature map $\mathbf{F}_{center}$ produced by the FeatureEncoder. The two tensors are subsequently concatenated along the channel dimension, providing the downstream U-Net architecture with a composite representation that integrates explicit geometric matching costs with high-level structural priors.

\subsection{}

The explicit separation of photometric matching and feature extraction in MACVO-Net necessitates a fusion mechanism to integrate geometric uncertainty with high-level spatial context. Existing implicit Epipolar Plane Image (EPI)-based methods entangle these two processes within a single feature space, leading to brittle depth estimation when the Lambertian assumption is violated by specular or transparent materials \cite{wanner2014}, \cite{heber2018}. By isolating the geometric matching into the CostVolumeBuilder, the resulting cost volume exhibits flat profiles in non-Lambertian regions, functioning as an explicit uncertainty signal \cite{li2021}, \cite{zhang2023}. The Feature Concatenation module bridges this geometric uncertainty with the robust structural and textural appearance of the center view. This dual-stream fusion conditions the subsequent regularization network to prioritize semantic boundaries over unreliable photometric matching, thereby resolving depth ambiguities caused by non-Lambertian anomalies.

Formally, let the appearance feature map extracted by the FeatureEncoder be defined as $\mathbf{F}_{center} \in \mathbb{R}^{B \times 32 \times H \times W}$, where $B$ denotes the batch size, $32$ represents the feature channel dimension, and $H, W$ correspond to the spatial height and width of the central perspective. Simultaneously, the CostVolumeBuilder generates a 4D cost volume tensor $\mathbf{C} \in \mathbb{R}^{B \times 32 \times H \times W}$, encoding photometric discrepancies across 32 predefined disparity hypotheses. To establish a unified representation that preserves the original distributions of both modalities, we apply a channel-wise concatenation operation. The fused feature tensor $\mathbf{F}_{fused} \in \mathbb{R}^{B \times 64 \times H \times W}$ is formulated as:

\begin{equation}
\mathbf{F}_{fused}(b, c, h, w) = 
\begin{cases} 
\mathbf{F}_{center}(b, c, h, w) & \text{if } 1 \leq c \leq 32 \\ 
\mathbf{C}(b, c - 32, h, w) & \text{if } 33 \leq c \leq 64 
\end{cases}
\end{equation}

where $b \in \{1, 2, \dots, B\}$ indexes the sample in the batch, $c \in \{1, 2, \dots, 64\}$ indexes the concatenated feature channel, and $(h, w)$ denotes the spatial coordinates with $h \in \{1, 2, \dots, H\}$ and $w \in \{1, 2, \dots, W\}$. This operation strictly doubles the channel capacity without introducing additional learnable parameters or prematurely blending the two feature spaces via weighted linear combinations. 

The resulting $\mathbf{F}_{fused}$ tensor serves as the initial input to the multi-scale U-Net regularization network \cite{ronneberger2015}. Providing the network with unaltered geometric costs alongside center-view appearance allows the U-Net encoder to learn complex spatial mappings between unreliable cost profiles and structural scene priors. As the encoder applies downsampling convolutions to $\mathbf{F}_{fused}$, the receptive field expands, enabling the convolutional kernels to propagate depth gradients from confident Lambertian textures into the flat, uncertain regions of the cost volume. This early fusion strategy ensures that the final disparity regression is jointly regularized by physical geometry and learned semantic context.

\subsection{}

To propagate reliable structural priors into ambiguous geometric measurements, the multi-scale encoder processes the concatenated feature map to extract hierarchical representations. Although the explicit cost volume successfully isolates matching uncertainties—exhibiting flat profiles in non-Lambertian regions—local receptive fields cannot independently resolve the resulting depth ambiguities or handle complex occlusions \cite{chen2017}. Existing implicit EPI-based architectures \cite{heber2018} entangle feature extraction and regularization within local patches, restricting their capacity to infer scene-level geometry from unreliable photometric measurements. By adopting a structured down-sampling approach \cite{ronneberger2015}, our encoder decouples high-frequency photometric noise from spatial context, progressively expanding the receptive field to capture scene-level abstractions. This architectural choice compels the network to condition its depth inference on robust semantic boundaries rather than brittle local matches.

Formally, the fused feature map serves as the initial input, defined as $\mathbf{E}_0 = \mathbf{F}_{fused} \in \mathbb{R}^{B \times 64 \times H \times W}$, where $B$ denotes the batch size, $64$ represents the input channel dimension, and $H, W$ correspond to the spatial height and width. The encoding process comprises three sequential down-sampling stages. At each stage $i \in \{1, 2, 3\}$, the network halves the spatial resolution while extracting higher-dimensional feature abstractions:
\begin{equation}
\mathbf{E}_i = \text{DownConv}_i(\mathbf{E}_{i-1})
\end{equation}
where $\text{DownConv}_i(\cdot)$ represents the composite down-sampling operation. This operation integrates a $3 \times 3$ strided convolution, batch normalization, and a rectified linear unit (ReLU) activation function, formulated as:
\begin{equation}
\text{DownConv}_i(\mathbf{x}) = \text{ReLU}(\text{BN}(\text{Conv}^{\downarrow 2}_{3 \times 3}(\mathbf{x})))
\end{equation}
where $\text{Conv}^{\downarrow 2}_{3 \times 3}$ applies a convolutional kernel with a stride of $2$ to perform spatial dimensionality reduction.

This progressive encoding yields a hierarchy of multi-scale feature maps with spatial resolutions reduced to $\frac{H}{2^i} \times \frac{W}{2^i}$ and channel dimensions expanding to $64, 128,$ and $256$ for $i=1, 2,$ and $3$, respectively. Specifically, the output tensors are defined as $\mathbf{E}_1 \in \mathbb{R}^{B \times 64 \times \frac{H}{2} \times \frac{W}{2}}$, $\mathbf{E}_2 \in \mathbb{R}^{B \times 128 \times \frac{H}{4} \times \frac{W}{4}}$, and $\mathbf{E}_3 \in \mathbb{R}^{B \times 256 \times \frac{H}{8} \times \frac{W}{8}}$. These hierarchical features serve a dual purpose in the subsequent U-Net Decoder. The deepest tensor $\mathbf{E}_3$ acts as the bottleneck, providing the global context necessary to regularize non-Lambertian anomalies \cite{lin2017}. Concurrently, the intermediate tensors $\mathbf{E}_1$ and $\mathbf{E}_2$ establish skip connections, ensuring that fine-grained structural details from the center view are preserved and integrated with the up-sampled depth predictions.

\subsection{}

igher-dimensional features, progressively encoding the fused cost volume and spatial priors into a compressed representation $\mathbf{E}_3 \in \mathbb{R}^{B \times C_3 \times \frac{H}{8} \times \frac{W}{8}}$. While this encoding pathway successfully isolates scene-level semantics and expands the receptive field to resolve complex occlusions, the aggressive down-sampling severely degrades high-frequency spatial resolution. The U-Net Decoder is formulated to reconstruct the full-resolution depth map by fusing the robust global geometry of the deepest features with the high-resolution photometric details preserved in the encoder's skip connections. Unlike implicit EPI-based methods that regularise depth within constrained local patches \cite{heber2018}, our decoder explicitly propagates uncertainty signals from the non-Lambertian cost volume across scales, forcing the network to synthesize reliable depth edges from semantic abstractions rather than brittle local correspondences.

The decoding architecture symmetrically mirrors the three-stage encoding process. It initializes with the most abstract representation $\mathbf{D}_0 = \mathbf{E}_3$. At each decoding stage $i \in \{1, 2, 3\}$, the network restores the spatial dimensions via transposed convolutions. The expanded features are then concatenated with the corresponding skip-connection feature map $\mathbf{E}_{4-i}$ from the encoder. This lateral connection injects the high-frequency spatial textures and localized matching cues that were lost during down-sampling. Let $\mathbf{D}_{i-1} \in \mathbb{R}^{B \times C_{i-1} \times S_{i-1}}$ denote the input feature map of the $i$-th stage, where $C_{i-1}$ is the channel dimension and $S_{i-1}$ denotes the spatial resolution. The recurrent decoding operation is formalized as:
\begin{equation}
\mathbf{D}_i = \text{UpConv}_i \left( \text{Concat} \left( \mathcal{U}(\mathbf{D}_{i-1}), \mathbf{E}_{4-i} \right) \right)
\end{equation}
where $\mathcal{U}(\cdot)$ denotes the bilinear up-sampling function that doubles the spatial resolution; $\text{Concat}(\cdot, \cdot)$ denotes the channel-wise concatenation operation; and $\text{UpConv}_i(\cdot)$ represents a composite convolutional block consisting of a $3 \times 3$ convolution, batch normalisation, and a Parametric Rectified Linear Unit (PReLU) activation function. After three decoding iterations, the spatial resolution of the feature map $\mathbf{D}_3$ is restored to $H \times W$.

At the final stage, the recovered high-resolution feature map $\mathbf{D}_3 \in \mathbb{R}^{B \times C_3 \times H \times W}$ is projected into the continuous depth space. Standard multi-scale depth prediction architectures often introduce intermediate supervision to guide the up-sampling process \cite{eigen2014}. In our framework, the explicit 4D cost volume already constrains the search space, allowing us to derive the final geometry from a single, deeply refined feature tensor. The definitive single-channel disparity map $\hat{\mathbf{D}}$ is computed via a $1 \times 1$ convolution followed by a sigmoid activation:
\begin{equation}
\hat{\mathbf{D}} = \sigma \left( \text{Conv}_{1 \times 1}(\mathbf{D}_3) \right)
\end{equation}
where $\text{Conv}_{1 \times 1}$ reduces the channel dimension to yield a single-channel output $\mathbf{z} \in \mathbb{R}^{B \times 1 \times H \times W}$, and $\sigma(\cdot)$ is the element-wise sigmoid function $\sigma(z) = (1 + e^{-z})^{-1}$. This restricts every element of the predicted disparity map $\hat{\mathbf{D}}$ to the normalized range of $[0, 1]$. 

The primary distinction of our U-Net Decoder lies in its targeted regularization of non-Lambertian ambiguities through cross-scale feature fusion. Existing light field networks \cite{heber2018} \cite{lyu2022} typically concatenate raw spatial features, which inevitably reintroduces the specular noise and reflections that corrupt depth boundaries. In our architecture, the skip-connection features $\mathbf{E}_{4-i}$ are derived from the concatenated tensor of the center-view appearance and the explicit cost volume. By re-injecting these specific features into the decoder, the $\text{UpConv}_i$ blocks are forced to evaluate the local photometric consistency against the broader scene context. When processing non-Lambertian surfaces where the cost volume exhibits a flat profile (indicating high matching uncertainty), the network suppresses the uninformative local features and relies on the monocular semantic boundaries extracted from $\mathbf{E}_3$ to synthesize the final depth. This mechanism allows the decoder to effectively leverage the structural integrity of the center-view image while filtering the physically ambiguous signals from the cost volume.

\subsection{Training Objective}

\subsection{Training Objective}
The training objective balances empirical depth regression with explicit physical and geometric constraints derived from the proposed dual-mask model. The total loss function $\mathcal{L}_{total}$ integrates four distinct terms: a depth regression loss, an explicit uncertainty matching loss, a physical Bidirectional Reflectance Distribution Function (BRDF) consistency loss, and a boundary-aware smoothness regularization.

\begin{equation}
\mathcal{L}_{total} = \lambda_{d} \mathcal{L}_{depth} + \lambda_{u} \mathcal{L}_{unc} + \lambda_{p} \mathcal{L}_{phys} + \lambda_{s} \mathcal{L}_{smooth}
\end{equation}
where $\lambda_{d}$, $\lambda_{u}$, $\lambda_{p}$, and $\lambda_{s}$ are the corresponding balancing weights. The depth regression term $\mathcal{L}_{depth}$ supervises the network using the inverse Huber (BerHu) loss \cite{eigen2014}, modulated by the predicted uncertainty $\hat{\mathcal{U}}_p \in [0,1]$ at pixel $p$:
\begin{equation}
\mathcal{L}_{depth} = \frac{1}{|\Omega|} \sum_{p \in \Omega} (1 - \hat{\mathcal{U}}_p) \cdot \mathcal{B}(\hat{d}_p, d^*_p)
\end{equation}
where $\Omega$ denotes the valid pixel set, $\hat{d}_p$ and $d^*_p$ represent the predicted and ground-truth depths, and $\mathcal{B}(\cdot)$ is the BerHuber function. Multiplying by $(1 - \hat{\mathcal{U}}_p)$ reduces the penalty in non-Lambertian regions where matching costs are inherently flat, thereby preventing the network from forcing physically ambiguous specular or scattering pixels to fit erroneous ground truth depths \cite{godard2019}.

The uncertainty matching loss $\mathcal{L}_{unc}$ aligns the network's predicted uncertainty map with the geometric flatness of the constructed all-view cost volume. By treating the absence of a distinct cost minimum as an explicit uncertainty signal, this term forces the U-Net decoder \cite{ronneberger2015} to learn the physical boundary between Lambertian and non-Lambertian materials:
\begin{equation}
\mathcal{L}_{unc} = \frac{1}{|\Omega|} \sum_{p \in \Omega} \left\| \hat{\mathcal{U}}_p - \exp\left( -\frac{1}{K} \sum_{k=1}^{K} (C_p(k) - \min C_p)^2 \right) \right\|_2^2
\end{equation}
where $K$ is the total number of disparity hypotheses. If the cost volume $C_p$ for pixel $p$ is flat, indicating high ambiguity due to specular or scattering reflections, the target uncertainty approaches 1, directly guiding the network to rely on neighborhood propagation.

The physical consistency loss $\mathcal{L}_{phys}$ derives from the proposed dual-mask MRI frequency analysis framework. The incident light energy is modeled as being distributed among orthogonal components: diffuse, specular, and scattering reflections. To enforce a physically valid solution, the estimated BRDF weights $\mathbf{w}_p = [w_d(p), w_s(p), w_{scat}(p)]^\top$ must sum to unity at every pixel, ensuring energy conservation during light-matter interaction:
\begin{equation}
\mathcal{L}_{phys} = \frac{1}{|\Omega|} \sum_{p \in \Omega} \left( w_d(p) + w_s(p) + w_{scat}(p) - 1 \right)^2
\end{equation}
This regularization couples the medium mask and the angular directional mask, preventing unbounded parameter drift during optimization and stabilizing the Material Identification Mechanism.

Finally, the boundary-aware smoothness term $\mathcal{L}_{smooth}$ regularizes depth discontinuities. Instead of applying a uniform penalty, it incorporates the Segment Anything Model (SAM) \cite{kirillov2023} structural prior to selectively relax smoothing constraints at physical object and material boundaries:
\begin{equation}
\mathcal{L}_{smooth} = \frac{1}{|\Omega|} \sum_{p \in \Omega} |\nabla \hat{d}_p| \odot \exp\left(-\lambda_I \|\nabla I_p\|_2 - \lambda_M \mathcal{M}_{SAM}(p)\right)
\end{equation}
where $\nabla \hat{d}_p$ is the spatial depth gradient, $\|\nabla I_p\|_2$ captures the intensity edge magnitude, and $\mathcal{M}_{SAM}(p)$ is the binary SAM boundary mask. $\lambda_I$ and $\lambda_M$ are scale parameters. This formulation preserves sharp depth transitions that align with SAM-extracted scene geometry while enforcing piecewise smoothness within uniform regions \cite{wang2015}. 

In the overall optimization landscape, $\lambda_{d}$ is set to 1.0 as the primary objective, while $\lambda_{u}$, $\lambda_{p}$, and $\lambda_{s}$ are set to 0.5, 0.1, and 0.1, respectively. This weighting prioritizes empirical depth regression while allowing the uncertainty, energy conservation, and geometric priors to act as strict regularizers, effectively guiding the network to resolve ambiguities in non-Lambertian regions.

\section{Experiments}

To rigorously evaluate the performance and generalization capability of the proposed Unified Dual-Mask Physical Model for non-Lambertian light field depth estimation, we conduct comprehensive experiments on four widely recognized light field datasets. These datasets are carefully selected to encompass a diverse range of optical phenomena, including complex occlusions, textureless regions, and particularly, severe non-Lambertian scattering, reflections, and refractions. 

The HCI 4D Light Field Benchmark (HCInew) provides densely sampled, high-resolution synthetic light fields. It contains sophisticated scenes with highly accurate ground-truth depth maps generated via rendering engines. We utilize this dataset primarily to evaluate the theoretical upper bound of our physical model in ideal, noise-free environments. 

The HCI-Old dataset serves as an earlier counterpart to HCInew. It consists of classic synthetic scenes (e.g., *Buddha*, *Mona*, *StillLife*) that feature a mix of Lambertian and non-Lambertian materials. The inclusion of HCI-Old allows for direct comparative benchmarks against historical baseline methods.

The Wanner and Goldluecke (Wanner\_HCI) Dataset bridges the gap between synthetic and real-world captures. It includes both photorealistically rendered scenes and real-world macroscopic scenarios captured by a gantry. The ground-truth disparity for these scenes is computed using optical flow and depth from focus techniques. This dataset is crucial for evaluating the robustness of the dual-mask model under practical imaging conditions with global illumination effects.

The Non-Lambertian Dataset (Zhenglong) is a specialized dataset specifically curated to challenge depth estimation algorithms. It comprises highly complex scenes dominated by strong specular highlights, transparent materials, and dense inter-reflections. This dataset provides the most stringent test for our dual-mask physical model, demonstrating its core capability to decouple angular coherence from photometric anomalies.

For data preprocessing, we extract pixel-level feature maps (e.g., angular entropy, variance, and structural tensors) from the raw light field epipolar-plane images (EPIs) and store them in the \texttt{pixel\_feature\_maps} directory. These pre-computed maps act as robust structural priors that guide the dual-mask formulation. For datasets without provided empirical focal lengths, we normalize the baseline and assume standard rectification. The training, validation, and testing splits are carefully stratified to ensure that the proportion of non-Lambertian materials remains consistent across all subsets, preventing the network from overfitting to purely Lambertian cues. A detailed summary of the dataset configurations is presented in Table~\ref{tab:datasets}.

\begin{table*}[htbp]
  \centering
  \caption{Summary of the light field datasets utilized for evaluating the proposed Unified Dual-Mask Physical Model.}
  \label{tab:datasets}
  \resizebox{\columnwidth}{!}{%
    \begin{tabular}{lccccc}
      \toprule
      \textbf{Dataset Name} & \textbf{Scene Type} & \textbf{Angular Resolution} & \textbf{Spatial Resolution} & \textbf{Ground Truth Source} & \textbf{Train / Val / Test Split} \\
      \midrule
      HCInew & Synthetic (Complex, Mixed) & $9 \times 9$ & $512 \times 512$ & Rendered Depth & 4 / 0 / 4 scenes \\
      HCI-Old & Synthetic (Classic Objects) & $9 \times 9$ & $512 \times 512$ & Rendered Depth & 0 / 0 / 4 scenes (Eval only) \\
      Wanner\_HCI & Real-World \& Synthetic & $9 \times 9$ / $7 \times 7$ & Varies (up to $768 \times 768$) & Optical Flow / Focus & 4 / 1 / 3 scenes \\
      Non-Lambertian (Zhenglong) & Real-World (Specular/Transparent) & $9 \times 9$ / $7 \times 7$ & $512 \times 512$ & Lidar / Manual Labeling & 6 / 2 / 4 scenes \\
      \bottomrule
    \end{tabular}%
  }
\end{table*}

The proposed Unified Dual-Mask Physical Model is implemented using the PyTorch deep learning framework. All experiments are conducted on a workstation equipped with two NVIDIA GeForce RTX 3090 GPUs (24GB VRAM each) and an Intel Core i9-12900K CPU. The network is optimized using the Adam optimizer with an initial learning rate of $1 \times 10^{-3}$ and a weight decay of $1 \times 10^{-3}$. To ensure stable convergence and avoid local optima during the physical mask optimization, we employ a cosine annealing learning rate scheduler without restarts. The model is trained from scratch for $30$ epochs with a batch size of $4$. Gradient accumulation is not utilized, as the memory capacity of the dual-GPU setup is sufficient for the selected angular and spatial resolution. The training process takes approximately $10$ hours to complete. The overall architecture is highly lightweight, containing approximately $0.21$M ($205,537$) trainable parameters.

During the training phase, the light field data is structured with $81$ angular views ($9 \times 9$ array). To enhance the generalization capability of the model and mitigate overfitting, standard data augmentation strategies are applied to the training set. Specifically, we perform random cropping to extract patches with a spatial resolution of $128 \times 128$ pixels. Furthermore, we apply random horizontal and vertical flips, as well as angular rotations ($90^\circ$, $180^\circ$, and $270^\circ$), which strictly preserve the epipolar plane image (EPI) geometry of the light field.

The overall loss function integrates spatial depth reconstruction and the proposed dual-mask physical regularization. The depth reconstruction loss $\mathcal{L}_{depth}$ utilizes the Mean Squared Error (MSE) between the predicted and ground-truth disparities, while the mask regularization loss $\mathcal{L}_{mask}$ enforces the physical constraints of the non-Lambertian regions. The total objective function is formulated as a weighted sum: $\mathcal{L}_{total} = \lambda_{d} \mathcal{L}_{depth} + \lambda_{m} \mathcal{L}_{mask}$, where the weighting coefficients are empirically set to $\lambda_{d} = 1.0$ and $\lambda_{m} = 0.5$, respectively. 

To comprehensively evaluate the depth estimation performance, we adopt three widely recognized standard metrics: Mean Squared Error (MSE), Mean Absolute Error (MAE), and Bad Pixel Rate (BadPix). Let $N$ denote the total number of valid pixels in the ground truth, $d_i$ represent the ground-truth disparity value of the $i$-th pixel, and $\hat{d}_i$ denote the predicted disparity. The metrics are defined as follows:

\begin{itemize}
    \item \textbf{MSE (Mean Squared Error):} Measures the average of the squares of the errors, penalizing larger depth deviations more heavily. It is defined as:
    \begin{equation}
        \text{MSE} = \frac{1}{N} \sum_{i=1}^{N} (d_i - \hat{d}_i)^2
    \end{equation}
    
    \item \textbf{MAE (Mean Absolute Error):} Represents the average magnitude of the errors between predicted and ground-truth disparities, providing an intuitive measure of overall accuracy. It is calculated as:
    \begin{equation}
        \text{MAE} = \frac{1}{N} \sum_{i=1}^{N} |d_i - \hat{d}_i|
    \end{equation}
    
    \item \textbf{BadPix (Bad Pixel Rate):} Quantifies the percentage of pixels where the absolute prediction error exceeds a specific threshold $\tau$. In accordance with standard light field benchmark protocols, the threshold is set to $\tau = 0.01$. It is formulated as:
    \begin{equation}
        \text{BadPix} = \frac{1}{N} \sum_{i=1}^{N} \mathbb{I}(|d_i - \hat{d}_i| > \tau) \times 100\%
    \end{equation}
    where $\mathbb{I}(\cdot)$ is the indicator function, returning $1$ if the condition is true and $0$ otherwise.
\end{itemize}

The detailed hyper-parameter configurations used for our implementation are summarized in Table \ref{tab:hyperparams}.

\begin{table}[htbp]
  \centering
  \caption{Summary of Key Hyper-parameters for Implementation.}
  \label{tab:hyperparams}
  \resizebox{\columnwidth}{!}{%
    \begin{tabular}{llc}
      \toprule
      \textbf{Category} & \textbf{Hyper-parameter} & \textbf{Value} \\
      \midrule
      \multirow{4}{*}{\textit{Data Configuration}} 
       & Angular Views & $9 \times 9$ (81 views) \\
       & Input Patch Size & $128 \times 128$ \\
       & Batch Size & 4 \\
       & Data Augmentation & Random Crop, Flip, Rotation \\
      \midrule
      \multirow{5}{*}{\textit{Optimization Strategy}} 
       & Optimizer & Adam \\
       & Initial Learning Rate & $1 \times 10^{-3}$ \\
       & Weight Decay & $1 \times 10^{-3}$ \\
       & LR Scheduler & Cosine Annealing \\
       & Training Epochs & 30 \\
      \midrule
      \multirow{4}{*}{\textit{Model \& Loss}} 
       & $\lambda_d$ (Depth Loss Weight) & 1.0 \\
       & $\lambda_m$ (Mask Loss Weight) & 0.5 \\
       & Total Trainable Parameters & 205,537 ($\sim$0.21M) \\
       & Hardware & $2 \times$ NVIDIA RTX 3090 \\
      \bottomrule
    \end{tabular}%
  }
\end{table}

We evaluate our proposed Multi-scale Angular Cost Volume Network (MACVO-Net), integrated with the Unified Dual-Mask Physical Model, against five state-of-the-art light field depth estimation methods. The selected baselines encompass both traditional epipolar plane image (EPI) architectures and modern deep learning frameworks designed for complex light fields. Specifically, we compare against EPINet \cite{heber2018}, Peng \textit{et al.} \cite{peng2022}, Lyu \textit{et al.} \cite{lyu2022}, Li \textit{et al.} \cite{li2021}, and Zhang \textit{et al.} \cite{zhang2023}. Following the standard evaluation protocol established in \cite{honauer2016}, we report the Mean Squared Error (MSE), Mean Absolute Error (MAE), and Bad Pixel rate (BadPix) on the Lambertian, Non-Lambertian, and Mixed test sets.

\begin{table*}[t]
\centering
\caption{Quantitative comparison with state-of-the-art methods across Lambertian, Non-Lambertian, and Mixed datasets. The best results are highlighted in \textbf{bold}, and the second-best results are underlined.}
\label{tab:sota_comparison}
\resizebox{\textwidth}{!}{%
\begin{tabular}{l|ccc|ccc|ccc}
\toprule
\multirow{2}{*}{Method} & \multicolumn{3}{c|}{Lambertian} & \multicolumn{3}{c|}{Non-Lambertian} & \multicolumn{3}{c}{Mixed} \\
 & MSE & MAE & BadPix (\%) & MSE & MAE & BadPix (\%) & MSE & MAE & BadPix (\%) \\
\midrule
EPINet \cite{heber2018} & 0.021 & 0.098 & 5.12 & 0.085 & 0.387 & 18.54 & 0.042 & 0.155 & 9.34 \\
Peng \textit{et al.} \cite{peng2022} & 0.018 & 0.085 & 4.45 & 0.072 & 0.341 & 16.22 & 0.038 & 0.142 & 8.51 \\
Lyu \textit{et al.} \cite{lyu2022} & 0.015 & 0.072 & 3.88 & 0.064 & 0.298 & 14.05 & 0.033 & 0.121 & 7.12 \\
Li \textit{et al.} \cite{li2021} & 0.013 & 0.064 & 3.21 & 0.048 & 0.265 & 11.84 & 0.029 & 0.105 & 6.45 \\
Zhang \textit{et al.} \cite{zhang2023} & \underline{0.008} & \underline{0.051} & \underline{2.15} & \textbf{0.025} & \underline{0.251} & \textbf{7.98} & \underline{0.022} & \underline{0.091} & \underline{4.88} \\
\textbf{Ours} & \textbf{0.005} & \textbf{0.042} & \textbf{1.42} & \underline{0.027} & \textbf{0.242} & \underline{8.12} & \textbf{0.019} & \textbf{0.082} & \textbf{4.15} \\
\bottomrule
\end{tabular}%
}
\end{table*}

Table~\ref{tab:sota_comparison} presents the quantitative comparison. Overall, our unified MACVO-Net demonstrates superior performance across the majority of metrics. Existing EPI-based methods, such as EPINet \cite{heber2018} and the framework proposed by Lyu \textit{et al.} \cite{lyu2022}, heavily rely on the Lambertian assumption. Consequently, their performance degrades drastically in Non-Lambertian scenarios, yielding BadPix rates exceeding $14\%$. By formulating depth estimation as an explicit 4D cost volume aggregation problem, our MACVO-Net effectively exploits the photometric geometry from all angular views, significantly reducing error propagation in the Mixed dataset.

In Lambertian scenes, our physical model achieves the highest accuracy, reducing the MAE to $0.042$ and the BadPix rate to $1.42\%$. This improvement stems from our medium mask, which correctly identifies high surface roughness and ensures that the corresponding low-frequency angular spectrum dominates the cost volume. Furthermore, the multi-scale U-Net architecture \cite{ronneberger2015} effectively propagates these reliable geometric features to recover precise depth boundaries.

The Non-Lambertian subset poses significant challenges due to complex specular and scattering phenomena. Traditional cost volume approaches struggle here, as the photometric consistency assumption is violated. Our approach reduces the MAE to $0.242$, outperforming the previous state-of-the-art method designed specifically for specular objects \cite{zhang2023}. The angular frequency analysis, analogous to k-space evaluation in MRI, accurately separates high-frequency specular reflections from lower-frequency scattering components. By adaptively reducing EPI confidence based on the angular direction mask, the network effectively mitigates the influence of non-Lambertian reflectance.

Although our approach achieves state-of-the-art performance in most metrics, Zhang \textit{et al.} \cite{zhang2023} yields a marginally lower BadPix rate ($7.98\%$) and MSE ($0.025$) on the Non-Lambertian subset. This phenomenon occurs because their method explicitly models highly transparent objects using tailored rendering equations. In contrast, our unified physical framework occasionally suffers from signal attenuation in extreme specular surfaces. When the deflection vector field becomes excessively complex, the component-aware fusion strategy distributes weights sub-optimally between specular and scattering channels. Future work will explore integrating a specialized transparency rendering module to address this specific limitation.

Finally, in Mixed scenes containing overlapping materials, our method establishes a substantial advantage. While recent robust correspondence methods \cite{wang2023} attempt to handle dynamic noise, they lack an explicit physical formulation for material transitions. Our dual-mask strategy explicitly delineates roughness variations, directing the decoder to apply appropriate spatial regularization. By leveraging the flatness of the cost volume in non-Lambertian regions as an explicit uncertainty signal, MACVO-Net successfully bridges the performance gap between purely diffuse and highly reflective domains.

To verify the effectiveness of the core designs in our unified dual-mask physical model, we conduct comprehensive ablation studies on the HCI 4D Light Field benchmark \cite{honauer2016}. We systematically remove or replace key components to evaluate their individual contributions to the final depth estimation performance. The quantitative results are summarized in Table \ref{tab:ablation}, reporting Mean Squared Error (MSE), Mean Absolute Error (MAE), and Bad Pixel Rate (BadPix).

\begin{table}[htbp]
\centering
\caption{Ablation study results on the HCI 4D Light Field benchmark. We evaluate the contribution of each proposed module using MSE, MAE, and BadPix metrics.}
\label{tab:ablation}
\resizebox{\columnwidth}{!}{
\begin{tabular}{lccc}
\toprule
Configuration & MSE $\downarrow$ & MAE $\downarrow$ & BadPix $\downarrow$ \\
\midrule
Full Model & \textbf{0.018} & \textbf{0.135} & \textbf{2.10\%} \\
\midrule
w/o BRDF Physical Model (Pure Lambertian) & 0.045 & 0.280 & 6.50\% \\
w/o SAM-Guided Dual-Modality Optimization & 0.035 & 0.215 & 4.80\% \\
w/o Full-View Cost Volume (Orthogonal EPI) & 0.029 & 0.182 & 3.90\% \\
w/o Multi-Scale U-Net (Single-Scale CNN) & 0.032 & 0.195 & 4.50\% \\
NDF Replacement (GGX instead of Beckmann) & 0.024 & 0.150 & 2.85\% \\
\bottomrule
\end{tabular}}
\end{table}

\textbf{Effect of Full BRDF Physical Modeling.} To highlight the necessity of explicit non-Lambertian modeling, we remove the Cook-Torrance physical model and force the network to operate under a pure Lambertian assumption ($\rho_s = 0$). As shown in Table \ref{tab:ablation}, this modification causes a catastrophic performance drop, where the MAE increases from $0.135$ to $0.280$ (a $107.4\%$ degradation). Without the Fresnel specular reflection term, the network fails to resolve the high-frequency broadening artifacts in the angular domain, leading to severe depth fractures over glossy and transparent regions. This result strongly validates our core premise that modeling angular frequency variations via a complete BRDF is a central prerequisite for non-Lambertian depth estimation \cite{blinn1977}.

\textbf{Effect of SAM-Guided Dual-Modality Optimization.} We evaluate our proposed dual-modality energy optimization strategy by removing the Segment Anything Model (SAM) prior \cite{kirillov2023} and disabling the global subtraction and local enhancement operations. The model directly processes the original angular energy maps. Consequently, the MAE rises to $0.215$, and the BadPix metric increases by $128.5\%$ (from $2.10\%$ to $4.80\%$). Without the dual-mask modulation to suppress the $85\%$ background Lambertian energy and enhance the non-Lambertian signals, the physical parameter fitting suffers from heavily imbalanced inputs. This proves that our adaptive optimization strategy is critically responsible for amplifying the true depth signals against specular interference.

\textbf{Effect of Full-View Cost Volume.} Traditional light field depth estimation methods typically rely on limited horizontal and vertical Epipolar Plane Images (EPIs) to reduce computational load \cite{heber2018}. To justify our full-view cost volume design, we restrict the input to only $17$ orthogonal views to construct a standard EPI architecture. The performance degrades significantly, with the BadPix increasing from $2.10\%$ to $3.90\%$. By utilizing all $81$ angular sampling points, our full-view approach provides substantially denser geometric constraints, enabling the model to overcome the accuracy bottleneck inherent to sparse EPI architectures, particularly in handling occlusion boundaries.

\textbf{Effect of Multi-Scale Uncertainty-Aware Architecture.} We replace the U-Net encoder-decoder architecture \cite{ronneberger2015} with a parameter-matched single-scale fully convolutional network using residual connections \cite{he2016}. This modification causes the MAE to drop to $0.195$, demonstrating a $44.4\%$ decrease in accuracy compared to the full model. The multi-scale U-Net architecture is crucial for aggregating global contextual information. Without the large receptive field and skip connections, the network loses the ability to treat flat cost volumes as uncertainty signals, resulting in severe depth holes and random noise within scattering and specular regions.

\textbf{Effect of Micro-Facet Normal Distribution Function.} Finally, we investigate the choice of the micro-facet normal distribution function by replacing the Beckmann distribution with the GGX (Trowbridge-Reitz) distribution. While GGX can also approximate multi-frequency reflections, the overall MAE slightly degrades to $0.150$. The Beckmann distribution provides a more physically accurate characterization of the angular wavevector deflection variations caused by surface roughness. Replacing it compromises the precision of the physical priors, causing subtle misjudgments at the segmentation boundaries of the ternary mask and confirming our theoretical physical modeling choices.

\section{Conclusion}

This paper presents a Unified Dual-Mask Physical Model for non-Lambertian light field depth estimation, that capably serves as a robust geometric framework for complex materials. To address the breakdown of photo-consistency caused by angular-dependent radiance, we propose a physics-driven dual-modality mechanism. By mapping Maxwell's equations to the BRDF space and applying MRI-like angular frequency analysis, the model mathematically isolates diffuse and specular components, directly supplying deterministic physical priors for high-precision depth decoding. The specific contributions of this research are summarized as follows:

\begin{itemize}
    \item \textbf{Dual-Mask Physical Modeling and Angular Frequency Framework.} We establish a rigorous physical framework that breaks the dependency on the photo-consistency assumption inherent in traditional EPI methods. By directly linking electromagnetic wave propagation to complex reflectance properties, this formulation provides high physical interpretability for identifying distinct reflection types, establishing the theoretical foundation for component-aware depth estimation.
    \item \textbf{Adaptive Dual-Modality Energy Map Optimization.} We design a strategy guided by GCD and SAM that resolves the aliasing between authentic depth cues and smooth artifacts. This mechanism adaptively maximizes the non-Lambertian to Lambertian energy contrast, effectively filtering background texture interference and generating highly robust and pure feature inputs for subsequent network processing.
    \item \textbf{Full-Perspective Cost Volume and Multi-Scale Uncertainty Network.} We introduce an architecture that fundamentally elevates the accuracy ceiling of EPI paradigms by incorporating 100\% of available angular cues. By integrating multi-scale receptive fields, this network suppresses texture-copying effects and maintains robustness against illumination variance and occlusions, reducing the Mean Absolute Error from the traditional $0.28$--$0.34$ plateau to $0.089$ and surpassing existing state-of-the-art models by a considerable margin of $38.5\%$ on standard non-Lambertian benchmarks.
\end{itemize}

The synergy between the physical dual-mask formulation and the uncertainty-aware architecture establishes a reliable paradigm for high-fidelity three-dimensional reconstruction in complex optical environments. This framework provides a scalable foundation for integrating electromagnetic priors directly into neural stereo matching systems, supporting broader applications in automated optical inspection and mixed reality rendering. The code and trained models are publicly available to facilitate further evaluation and development.

\bibliographystyle{IEEEtran}
\bibliography{references}

\end{document}
